% Generated draft: requires mathematical and source review.
\documentclass[11pt]{article}
\usepackage[margin=25mm]{geometry}
\usepackage{fontspec}
\setmainfont{Latin Modern Roman}
\usepackage{amsmath,amssymb,mathtools,unicode-math}
\setmathfont{Latin Modern Math}
\usepackage{enumitem,needspace,xurl,hyperref}
\setlength{\parindent}{0pt}
\setlength{\parskip}{5pt}
\setlength{\emergencystretch}{3em}
\newcommand{\N}{\mathbb{N}}
\newcommand{\Z}{\mathbb{Z}}
\newcommand{\Q}{\mathbb{Q}}
\newcommand{\R}{\mathbb{R}}
\newcommand{\C}{\mathbb{C}}
\newcommand{\F}{\mathbb{F}}
\newcommand{\A}{\mathbb{A}}
\newcommand{\PP}{\mathbb{P}}
\newcommand{\E}{\mathbb{E}}
\newcommand{\eps}{\varepsilon}
\newcommand{\dd}{\,\mathrm{d}}
\begin{document}
\section*{UnsolvedMath ID: 30003070}
\textbf{Source problem code:} catalogue entry
\subsection{Definitions and mathematical statement}
A plabic graph is a planar bicolored graph embedded in a disk, with designated boundary vertices and black or white internal vertices. The supplied record does not specify which plabic graphs are admissible, which Grassmannian they represent, or the construction assigning a valuation to such a graph.

For clarity, the usual Newton--Okounkov body construction can be described after making the following choices. Let \(X\) be an integral projective variety of dimension \(d\) over a field \(K\), let \(L\) be a line bundle defining the relevant polarization, let \(\tau\) be a nonzero rational section of \(L\), and let
\[
 v:K(X)^\times\longrightarrow (\Z^d,\preceq)
\]
be a valuation, where \(\preceq\) is a total order compatible with addition. Thus \(v(fg)=v(f)+v(g)\) and \(v(f+g)\succeq\min\{v(f),v(g)\}\) whenever the expressions are defined. With these choices, write
\[
 \Delta_v(X,L;\tau)
 =\overline{\operatorname{conv}}\!\left(
 \bigcup_{m\geq 1}
 \left\{\frac{1}{m}v(s/\tau^m):
 0\neq s\in H^0(X,L^{\otimes m})\right\}
 \right)\subseteq\R^d.
\]
The field, polarization, valuation conventions, and normalization in this display are explanatory parameters; the source record does not fix them.

In the usual root-subgroup interpretation, a birational sequence for a Grassmannian \(X\) is an ordered sequence \(S\) of one-parameter root subgroups whose product, followed by passage to the Grassmannian, gives a birational parametrization \(\A_K^d\dashrightarrow X\). An admissible total order on exponent vectors is then used to select lowest terms and construct a valuation, denoted here by \(v_{S,\prec}\). The precise admissibility conditions on \(S\) and \(\prec\) are not supplied. In particular, the phrase ``proper total order'' in the supplied literature assessment is not defined there.

\textbf{Question, with unresolved conventions explicit.} Does every Newton--Okounkov body arising from a plabic graph admit a realization by a birational sequence for the relevant Grassmannian together with a suitable total order? Schematically, after fixing the intended graph construction, polarization, normalization, and meaning of realization, is it true that
\[
 \forall\,\Gamma\text{ admissible},\qquad
 \exists\,S\text{ a birational sequence},\quad
 \exists\,\prec\text{ an admissible total order},\qquad
 \Delta_{\Gamma}\sim\Delta_{S,\prec}?
\]
Here \(\Delta_{\Gamma}\) denotes the body supplied by the intended plabic-graph construction, and \(\Delta_{S,\prec}\) denotes the body supplied by the intended birational-sequence construction. The relation \(\sim\) is deliberately unresolved: the record does not say whether ``obtained'' requires equality in identified coordinates, integral affine equivalence, or another specified equivalence. Nor does it explicitly settle whether both constructions must use the same Grassmannian and polarization. These omissions prevent this schematic question from being a uniquely specified quantified assertion.
\subsection{Short English statement}
The question asks whether every Newton--Okounkov body constructed from a plabic graph can also be constructed using a birational sequence for a Grassmannian and a suitable total order. The graph class, geometric data, allowable orders, and meaning of agreement between bodies require clarification. The supplied record labels the problem open; this draft does not independently verify that status.
\subsection{Sources}
\begin{itemize}
\item Supplied local UnsolvedMath registry record, permanent ID \(30003070\), source problem code \texttt{catalogue entry}, titled ``Birational Sequences for Plabic Newton--Okounkov Bodies,'' in the collection ``Oberwolfach Reports Open Problems.'' The record gives proposed year 2016 and names no individual proposer.
\item Supplied source citation: \emph{Mini-Workshop: PBW Structures in Representation Theory} (2016), Oberwolfach Reports, pp.~615--651, DOI \texttt{10.4171/owr/2016/13}. \url{https://doi.org/10.4171/owr/2016/13}. The record describes Oberwolfach Reports as documenting mathematical workshops and their open-problem discussions; no exact problem page is supplied.
\item Supplied literature source labeled ``birational sequences'': \url{https://arxiv.org/abs/1510.02295}.
\item Supplied literature source labeled ``combinatorial mutations'': \url{https://arxiv.org/abs/2107.04264}.
\item Literature assessment dated exactly \texttt{2026-08-21}: the supplied curation labels the problem open and reports that birational sequences recover many plabic-graph Newton--Okounkov bodies, while no general theorem establishes the proposed realization for every such body with a proper total order. This is supplied dated curation, not an independently verified literature assessment. The record expressly requests independent verification of time-sensitive status and cited claims.
\item The supplied registry labels the statement \texttt{corrected\_verified} and reports a manual semantic, typographic, and LaTeX review. These are source metadata, not mathematical verification performed for this draft. None of the linked sources was consulted in preparing this response.
\end{itemize}
\end{document}
